% =============================================================================
% 第2章：数据来源与探索性分析
% =============================================================================

\section{数据来源与探索性分析}
\label{sec:data}

\subsection{数据来源与金融匿名化背景}

本文采用Kaggle平台上的Credit Card Fraud Detection Dataset 2023\cite{kaggle2023fraud}。该数据集包含欧洲信用卡持卡人的真实交易记录，原始数据共568,630条交易，31个特征列：一个交易ID列（id）、28个经PCA匿名化处理的特征（V1--V28）、交易金额（Amount）和类别标签（Class，0=正常交易，1=欺诈交易）。

PCA匿名化在此数据集中的使用并非学术实验的人为设计，而是金融行业数据共享实践的缩影。出于客户隐私保护和监管合规要求，金融机构在对外发布或跨机构共享交易数据时，必须对涉及个人身份和财务行为信息的原始特征进行脱敏。PCA通过将原始特征投影到一组彼此正交的主成分上，在保留数据全局方差结构的同时，使得攻击者无法直接还原原始特征的具体含义或反推特定客户的身份信息。正因如此，Kaggle等公开数据平台上的金融数据集普遍采用PCA匿名化处理——这一"行业惯例"恰好构成了本文研究问题的现实基础。需要特别指出，本文并不批评PCA匿名化本身——从隐私保护角度它是合理且必要的——而是诊断在其之上叠加深度学习模型时产生的意外后果。

数据集的类别平衡处理如下。原始数据中正负类接近1:1平衡（欧洲信用卡交易的真实欺诈率通常在0.1\%--1\%之间）。为模拟真实场景中约0.5\%的欺诈发生率，本文对多数类（正常交易）进行随机下采样并限制最大样本数为150,000条，同时从少数类（欺诈交易）中随机抽取750条，使最终欺诈比例约为0.5\%（正负比约1:200）。这一处理对于本研究的实验目的——在极度不平衡条件下诊断自注意力的行为——是必要的。

\subsection{数据集特征}

经处理后的数据集包含150,750条记录（150,000正常 + 750欺诈），29个数值型特征（V1--V28 + Amount），无缺失值。

图\ref{fig:class_balance}呈现了经不平衡处理后的完整数据集类别分布：正常交易样本数为150,000，欺诈交易样本数仅为750，正负类比例约为200:1。这种极端的类别不平衡是信用卡欺诈检测任务的核心挑战——模型在训练过程中极易被多数类主导，导致对少数类的识别能力严重退化。仅0.5\%的欺诈比例意味着，即使一个"将所有交易预测为正常"的朴素分类器也能达到99.5\%的准确率——这正是准确率（Accuracy）在不平衡场景下完全失效的原因，也解释了本文为何采用AUPRC而非准确率作为主要评估指标。

\begin{figure}[htbp]
    \centering
    \includegraphics[width=0.72\textwidth]{class_balance.pdf}
    \caption{经不平衡处理后的数据集类别分布。正常交易150,000条，欺诈交易750条，正负比约200:1，准确率在此场景下完全失效。}
    \label{fig:class_balance}
\end{figure}

图\ref{fig:amount_dist}以核密度估计对比了两类交易的金额分布。蓝色曲线（正常交易，$n{=}150{,}000$）和橙色曲线（欺诈交易，$n{=}750$）均呈现明显的右偏（长尾）分布，但存在两个关键差异：第一，欺诈交易的KDE峰值位置略高于正常交易，表明欺诈者倾向于选择中等偏高的交易金额以在"隐匿性"与"收益"之间取得平衡；第二，在高金额区间（$>15,000$标准化单位），欺诈交易的密度相对更高，这与银行业务经验一致——高价值交易是欺诈者的优先目标。然而，两类分布在整体上存在大量重叠区域，意味着仅凭交易金额远不足以区分欺诈与正常交易，必须依赖多维特征联合判别。

\begin{figure}[htbp]
    \centering
    \includegraphics[width=0.75\textwidth]{amount_distribution.pdf}
    \caption{正常交易与欺诈交易的金额核密度估计曲线。蓝色为正常交易，橙色为欺诈交易，两类均呈右偏分布且存在大量重叠，欺诈交易在高金额区间密度略高。}
    \label{fig:amount_dist}
\end{figure}

\subsection{数据预处理}

预处理流水线包含以下步骤：

\begin{enumerate}
    \item \textbf{特征选择}：移除ID列（不携带预测信息），保留28个PCA特征和Amount共29个特征。
    \item \textbf{异常值处理}：对Amount特征采用IQR方法进行异常值截断（$Q1 - 1.5 \times IQR$ 至 $Q3 + 1.5 \times IQR$），仅在训练集上计算上下界以避免数据泄露。V1--V28为PCA变换后的特征，其分布为原始特征的线性组合，不具有传统意义上的"异常值"含义，因此不做截断处理。
    \item \textbf{标准化}：使用StandardScaler对所有29个数值特征进行标准化（均值为0，标准差为1），仅在训练集上拟合参数，应用于验证集和测试集。
    \item \textbf{数据集划分}：按70\%/15\%/15\%比例进行分层划分（Stratified Split），确保各子集中欺诈比例一致。最终产出：训练集105,525条，验证集22,613条，测试集22,612条。
\end{enumerate}

上述EDA和预处理结果表明：（1）任何有效的检测模型必须采用针对性的不平衡处理策略，因为0.5\%的欺诈率使标准分类器失效；（2）多特征联合判别是必需的，因为单一金额特征区分度不足；（3）评估时必须使用对类别不平衡敏感的指标（AUPRC而非准确率或ROC-AUC）；（4）所有V1--V28特征为PCA匿名化后的成分——这一事实贯穿全文，是注意力塌缩发生的结构性前提。
